\documentclass[11pt,a4paper]{article}

% --- packages ---
\usepackage[T1]{fontenc}
\usepackage[utf8]{inputenc}
\usepackage{lmodern}
\usepackage[margin=1in]{geometry}
\usepackage{microtype}
\usepackage{xcolor}
\usepackage{titlesec}
\usepackage{enumitem}
\usepackage{parskip}
\usepackage{hyperref}
\usepackage{tcolorbox}
\usepackage{fancyhdr}
\usepackage{setspace}
\usepackage{booktabs}

% --- colors ---
\definecolor{accent}{HTML}{1F3A93}
\definecolor{soft}{HTML}{4A4A4A}
\definecolor{quotebg}{HTML}{F4F1EC}
\definecolor{quoteborder}{HTML}{C8B273}
\definecolor{cuebg}{HTML}{ECF1F8}

% --- hyperref ---
\hypersetup{
  colorlinks=true,
  linkcolor=accent,
  urlcolor=accent,
  pdftitle={Nirnay --- Hackathon Demo Video Script},
  pdfauthor={Roshan Yadav}
}

% --- section styling ---
\titleformat{\section}
  {\color{accent}\Large\bfseries}
  {\thesection.}{0.6em}{}
\titleformat{\subsection}
  {\color{accent}\large\bfseries}
  {}{0em}{}

% --- header / footer ---
\pagestyle{fancy}
\fancyhf{}
\renewcommand{\headrulewidth}{0pt}
\fancyhead[L]{\small\color{soft}\textit{Nirnay --- Demo Video Script}}
\fancyhead[R]{\small\color{soft}Roshan Yadav}
\fancyfoot[C]{\small\color{soft}Page \thepage}

% --- spoken-line block ---
\newtcolorbox{spoken}{
  enhanced,
  colback=quotebg,
  colframe=quoteborder,
  boxrule=0pt,
  leftrule=3pt,
  arc=2pt,
  left=10pt, right=10pt, top=8pt, bottom=8pt,
  fontupper=\itshape
}

% --- visual cue block ---
\newtcolorbox{cue}{
  enhanced,
  colback=cuebg,
  colframe=cuebg,
  boxrule=0pt,
  arc=2pt,
  left=10pt, right=10pt, top=6pt, bottom=6pt,
  fontupper=\small
}

% --- timestamp tag ---
\newcommand{\ts}[1]{{\small\color{accent}\texttt{[#1]}}}

\begin{document}

% --- title block ---
\begin{center}
  {\Huge\bfseries\color{accent} Nirnay\,---\,Demo Video Script}\\[6pt]
  {\large\itshape Every verdict, with evidence.}\\[4pt]
  {\small Theme 3\,$\cdot$\,CRPF Procurement \quad|\quad PanIIT AI for Bharat Hackathon 2026}\\[2pt]
  {\small Presenter: \textbf{Roshan Yadav} \quad|\quad Target runtime: \textbf{2\,min 30\,s}}
\end{center}

\vspace{6pt}
\hrule height 0.4pt
\vspace{14pt}

\noindent\textbf{How to read this script.}
Italic blocks on a parchment background are the lines you speak --- straight to camera, conversational, with the small pauses written in. Light-blue blocks are visual cues: what's on screen, what to click, where to look. Timestamps are guides, not handcuffs --- if a beat lands, give it room.

\vspace{4pt}

% =====================================================
\section{Opening --- on camera, looking at lens}
\ts{0:00 -- 0:15}

\begin{spoken}
``Hi, I'm Roshan.

\smallskip
Before I show you what I built, let me ask you something. Imagine you're a CRPF procurement officer. A tender lands on your desk\,---\,a hundred-page PDF, fifteen bidders, each with eight to ten supporting documents. Your job? Decide who's eligible. And defend that decision \emph{in court}, sometimes years later.''
\end{spoken}

\begin{cue}
\textbf{Direction.} Pause here. Let it breathe. Slight shift in tone before the next line.
\end{cue}

\begin{spoken}
``Right now, that takes officers about \textbf{two weeks} per tender. I built something that does it in under a minute\,---\,and unlike every other AI tool out there, \textbf{mine can actually be defended in court.} It's called \textbf{Nirnay}\,---\,Hindi for \emph{verdict}. And it's built specifically for Theme 3.''
\end{spoken}

% =====================================================
\section{The problem}
\ts{0:15 -- 0:35}

\begin{cue}
\textbf{Cut to.} Screen recording starts. Briefly: B-roll of a desk with scattered PDFs, then transition into the dashboard.
\end{cue}

\begin{spoken}
``Here's what makes this hard. Existing AI tools give you a black-box yes-or-no. An officer can't walk into a courtroom and say \emph{`the AI said so.'} That's not a defence\,---\,that's a liability.

\smallskip
The other option? Build an OCR pipeline with LayoutLM, PaddleOCR, six months of fine-tuning. Defence Ministry timelines don't work that way.

\smallskip
So I asked one question: \textbf{what would a system look like that an officer could actually trust?} Two non-negotiables fell out of that question.''
\end{spoken}

% =====================================================
\section{The two principles}
\ts{0:35 -- 0:55}

\begin{cue}
\textbf{On screen.} Clean text overlay, two principles side by side. Simple typography, lots of whitespace.
\end{cue}

\begin{spoken}
``\textbf{One\,---\,every verdict points at evidence.} Not a number floating in a dashboard. The exact quote, the exact document, the exact page number. If the system says a bidder's turnover is five crore, you can click it and see the line on the CA certificate that proves it.

\smallskip
\textbf{Two\,---\,Nirnay never silently rejects.} If evidence is missing or unclear, the verdict is \texttt{needs\_review}. Never \texttt{not\_eligible}. Because rejecting the wrong bidder costs the government legal cases and reputation. A human review costs ten minutes.''
\end{spoken}

% =====================================================
\section{The demo --- live walkthrough}
\ts{0:55 -- 1:45}

\begin{cue}
\textbf{Cut to app.} Dashboard view. Move slowly --- the click-to-evidence moment is the hero of the whole pitch. Don't rush.
\end{cue}

\begin{spoken}
``Okay, let me show you. \emph{(Open dashboard.)} This is a real tender. Five bidders, ten eligibility criteria already extracted by Gemini 2.5 Flash\,---\,directly from the PDF, no OCR pipeline.

\smallskip
\emph{(Click Evaluation page.)} I hit \textbf{Run Evaluation}. Watch this\,---\,\textbf{two seconds}, all five bidders evaluated against all ten criteria. That's fifty verdicts.

\smallskip
\emph{(Click on Sharma Construction\,$\rightarrow$\,eligible.)} Eligible. But here's the thing\,---\,\emph{(click a verdict)}\,---\,every single verdict opens an evidence panel. Look. The exact quoted line from their CA certificate. Document name. Page number. Confidence score. Reasoning. \textbf{An officer can take a screenshot of this and put it on file.}

\smallskip
\emph{(Click Gupta Builders\,$\rightarrow$\,not eligible.)} This one's rejected. Why? Their turnover is 3.4 crore. The tender requires 5. Look\,---\,\emph{(point to evidence)}\,---\,the source line is right there. Defensible.

\smallskip
\emph{(Click National Infrastructure\,$\rightarrow$\,needs review.)} And this is the one I'm proudest of. The system \emph{doesn't know.} Their ISO certificate expires \emph{on} the bid date. Is that valid or not? The rules are ambiguous. So Nirnay refuses to guess\,---\,it flags it for an officer to decide. \textbf{That's the rule: when in doubt, ask a human.}''
\end{spoken}

% =====================================================
\section{The audit chain --- the killer demo}
\ts{1:45 -- 2:10}

\begin{cue}
\textbf{Hero shot.} Pause for a beat after ``Caught.'' Let the red row do the work. This is the moment the pitch lands.
\end{cue}

\begin{spoken}
``Now here's the part I think judges will care about most.

\smallskip
\emph{(Open Audit page.)} Every action in the system\,---\,every upload, every verdict, every override\,---\,goes into a hash-chained audit log. SHA-256, computed in the browser. Sixty-plus events here.

\smallskip
\emph{(Click `Simulate tampering'.)} Watch\,---\,I'm going to corrupt one row. Directly via SQL. Like an attacker would.

\smallskip
\emph{(Click `Verify Chain'.)} One click. \textbf{Caught.} Red row. Exact event flagged. \textbf{This is what makes a verdict survive a court challenge five years from now.}''
\end{spoken}

% =====================================================
\section{Closing --- back on camera}
\ts{2:10 -- 2:30}

\begin{spoken}
``So that's Nirnay.

\smallskip
One officer, one tender\,---\,\textbf{two weeks of work, done in minutes.} A hundred percent citation coverage. A tamper-proof audit trail. Built on free-tier Gemini, deployed as a static site, costs less than fifty rupees per tender to run.

\smallskip
If a CRPF directorate processes a hundred and twenty tenders a year, that's \textbf{twelve thousand four hundred officer-hours} redirected away from paperwork\,---\,and into the work that actually needs human judgement.

\smallskip
I'm Roshan Yadav. This is Nirnay. \textbf{Every verdict, with evidence.}

\smallskip
Thank you.''
\end{spoken}

% =====================================================
\section{Delivery notes for recording}

\begin{itemize}[leftmargin=*, itemsep=4pt]
  \item \textbf{Pace yourself on the demo segment.} The temptation is to rush. Don't. The ``click and see evidence'' moment is the entire pitch in one beat\,---\,let it land.
  \item \textbf{The tampering demo is your hero shot.} Pause for half a second after ``Caught.'' The red row is doing the talking; you don't have to.
  \item \textbf{Speak like you're explaining it to a friend who's smart but tired.} Not a panel. Not a marketing video. Keep contractions ($I'm$, $doesn't$); keep the small pauses already written in.
  \item \textbf{Theme 3 callout} sits in the opening line on purpose. Don't bury it later.
  \item \textbf{If you go over time}, the easiest cut is the closing economics paragraph (the 12{,}400 officer-hours line). It's nice-to-have, not load-bearing.
  \item \textbf{Camera setup.} Eye-level. Soft front light, no harsh shadow. Plain, slightly out-of-focus background\,---\,not a blank wall.
  \item \textbf{Audio matters more than video.} If you can borrow a lavalier or a decent USB mic, do it. Phone audio reads as ``hackathon submission''; clean audio reads as ``founder.''
\end{itemize}

\vspace{10pt}

% =====================================================
\section*{One-line summary}
\addcontentsline{toc}{section}{One-line summary}

\begin{tcolorbox}[
  enhanced, colback=accent!8, colframe=accent, boxrule=0.6pt, arc=3pt,
  left=12pt, right=12pt, top=10pt, bottom=10pt
]
\centering\large\itshape
``Two weeks of paperwork. One minute of AI. Every verdict pointed at the exact line of evidence.\\
That's Nirnay.''
\end{tcolorbox}

\end{document}
